%% annex/p07/normalization.tex --- source-status-preserving P7.3 sketches.

\section{ND and typed-term normalization}

\begin{annexitem}{annex:p7.norm.1}{thm:p07-normalization}{Normalization for intuitionistic natural deduction}
\sketchnote{the complete principal-reduction table for
\(\land,\lor,\to,\forall,\exists\); the complete permutation-conversion
table for \(\lor E\) and \(\exists E\); preservation of every discharge and
eigenvariable condition; and the proof that the maximal-formula-rank and
maximal-segment-length measure decreases in every case.}
Principal reductions remove an introduction--elimination detour; permutation
conversions move a later elimination above a disjunction or existential
elimination.  Alpha-renaming preserves capture avoidance and the registered
eigenvariable conditions.  Order derivations lexicographically by the
greatest rank of a maximal formula and the lengths of maximal segments.
Completion of the omitted case tables would show that repeated conversion
terminates in a normal derivation with the same conclusion and no new open
assumption formula.
\end{annexitem}

\begin{annexitem}{annex:p7.norm.2}{thm:p07-stlc-subject-reduction}{Subject reduction for STLC}
\sketchnote{the full induction proving the typing-substitution lemma,
including variable-renaming and capture-avoidance cases, and the induction
over every term-context clause in the compatible closure of beta reduction.}
First establish typing substitution: from
\(\Theta_{\mathrm{ty}},x{:}\beta
 \derives_{\mathsf{STLC}}M{:}\alpha\) and
\(\Theta_{\mathrm{ty}}\derives_{\mathsf{STLC}}N{:}\beta\), infer
\(\Theta_{\mathrm{ty}}\derives_{\mathsf{STLC}}M[N/x]{:}\alpha\).
Apply it to an outermost beta contraction.  An induction over the compatible
term context then propagates the preserved type when the contracted redex
occurs inside an abstraction or either side of an application.
\end{annexitem}

\begin{annexitem}{annex:p7.norm.3}{thm:p07-stlc-strong-normalization}{Strong normalization for STLC}
\sketchnote{the neutral-term and beta-expansion closure lemmas for the
reducibility candidates, their stability under reduction and substitution,
and every case of the fundamental lemma for the registered typing rules.}
Interpret a base type by the strongly normalizing terms.  Interpret
\(\alpha\mathbin{\to_{\mathrm{ty}}}\beta\) by the terms whose application to
every candidate member at \(\alpha\) lies in the candidate at \(\beta\).
The fundamental lemma, proved by induction on a typing derivation and a
candidate-valued substitution for its context, places every well-typed term
in its type's candidate.  The candidate at every type consists of strongly
normalizing terms, yielding the stated conclusion once the omitted closure
lemmas are supplied.
\end{annexitem}
